
\chapter{Análisis de hallazgos y estrategias de mitigación}

\section{Priorización de hallazgos}

Los hallazgos provienen de dos fuentes: las herramientas de seguridad
(Semgrep, Trivy, Gitleaks, TruffleHog y OWASP ZAP) y el modelado de amenazas
STRIDE. Para ordenarlos se consideraron tres criterios:

\begin{enumerate}
    \item \textbf{Severidad:} la que reporta la herramienta (por ejemplo,
    alta para el CVE de Trivy y para la versión de Nginx en ZAP) o, en las
    amenazas de STRIDE, el riesgo calculado como impacto por probabilidad.
    \item \textbf{Impacto:} qué se pierde si el hallazgo se explota
    (confidencialidad, integridad o disponibilidad de los datos de
    \texttt{moon\_notes}).
    \item \textbf{Criticidad del componente y facilidad de explotación:} se
    prioriza lo que es alcanzable desde la red sin credenciales, como la API,
    por encima de lo que exige acceso previo al servidor.
\end{enumerate}

Cuando un hallazgo de una herramienta coincide con una amenaza del modelo, se
trata como una sola entrada y se indica su origen. Esto ocurre con el CORS
abierto, las cabeceras ausentes y los errores expuestos, que el análisis
STRIDE predijo y ZAP confirmó. Los 15 avisos de Semgrep comparten una misma
regla y una misma causa, por lo que se agrupan en un único hallazgo.

\section{Tabla de hallazgos identificados}

\begin{longtable}{
    p{0.8cm}
    p{4.1cm}
    p{1.4cm}
    p{2.3cm}
    p{5.0cm}
}

\caption{Hallazgos de seguridad y recomendaciones de mitigación.}
\label{tab:hallazgos}\\

\toprule
\textbf{ID} &
\textbf{Descripción} &
\textbf{Severidad} &
\textbf{Herramienta} &
\textbf{Recomendación de mitigación} \\
\midrule
\endfirsthead

\multicolumn{5}{c}{
    \tablename\ \thetable{} -- continuación
}\\
\toprule
\textbf{ID} &
\textbf{Descripción} &
\textbf{Severidad} &
\textbf{Herramienta} &
\textbf{Recomendación de mitigación} \\
\midrule
\endhead

\midrule
\multicolumn{5}{r}{Continúa en la siguiente página}\\
\endfoot

\bottomrule
\endlastfoot

001 & La API no tiene autenticación ni autorización: cualquiera puede
leer, crear, modificar y borrar notas (S1, T1). &
Alta & Modelado STRIDE &
Añadir autenticación (sesión o JWT) y autorización por recurso. Como
mínimo, proteger \texttt{POST} y \texttt{DELETE} con un token. \\

002 & Dependencia \texttt{body-parser} 1.20.8 con CVE-2026-111181:
denegación de servicio por error no controlado en la descompresión. &
Alta & Trivy &
Actualizar \texttt{body-parser} a 1.20.9 en \texttt{backend} y regenerar
\texttt{package-lock.json}. \\

003 & Nginx expone una versión con vulnerabilidades conocidas en la
cabecera \texttt{Server}. &
Alta & OWASP ZAP &
Actualizar la imagen de Nginx a una versión corregida, fijarla por
\textit{digest} y ocultar la versión (\texttt{server\_tokens off}). \\

004 & Faltan cabeceras de seguridad (CSP, anti-clickjacking,
\texttt{X-Content-Type-Options}) y no hay TLS (T4). &
Media & OWASP ZAP &
Añadir las cabeceras con \texttt{add\_header} en Nginx y terminar TLS en el
despliegue. \\

005 & CORS abierto: la API responde
\texttt{Access-Control-Allow-Origin: *} (T3). &
Media & OWASP ZAP &
Restringir el origen con \texttt{cors(\{ origin: <frontend> \})}. \\

006 & Acciones de GitHub referenciadas por etiqueta mutable en el pipeline
(15 avisos de una misma regla). &
Media & Semgrep &
Fijar cada acción por el \textit{hash} completo del commit. \\

007 & Fuga de información: errores 500 con \texttt{err.message} y
cabeceras \texttt{X-Powered-By} y \texttt{Server} (I1). &
Baja & OWASP ZAP y STRIDE &
Responder con mensajes genéricos, registrar el detalle en el servidor y
desactivar \texttt{X-Powered-By}. \\

\end{longtable}

Gitleaks y TruffleHog no generaron hallazgos, por lo que no figuran en la
tabla. Tampoco se incluyen aquí las amenazas del modelo de riesgo menor o no
verificadas (candidatas), que se describen en el capítulo de modelado de
amenazas. Entre las confirmadas solo por el análisis del código, y no por
herramientas, destacan la ausencia de límites de tasa (D1) y el puerto 3001
publicado (S2).

\section{Análisis de las mitigaciones}

\subsection{Estado actual}

En el momento de redactar este informe \textbf{ninguno de los siete
hallazgos de la tabla ha sido corregido}. El trabajo realizado se centró en
construir el modelo de amenazas, implementar las herramientas de análisis y
centralizar sus resultados en DefectDojo. Los hallazgos están por tanto
\textbf{pendientes}, y la corrección se plantea como la siguiente fase.

Sí se aplicaron mejoras que no corrigen un hallazgo concreto pero reducen la
superficie de configuración: los valores de \texttt{compose.yaml} y de los
\texttt{Dockerfile} (imágenes, puertos y variables) dejaron de estar escritos
a mano y ahora se leen de variables del repositorio. Además, se mantienen
vigentes los controles existentes: consultas SQL parametrizadas, PostgreSQL
en una red interna y sin publicar, contenedores sin permisos de
administrador, contraseña como \textit{Docker secret} fuera de Git y rama
\texttt{main} protegida con revisión obligatoria.

\subsection{Plan de corrección propuesto}

Las acciones se ordenan por esfuerzo e impacto, de modo que las más baratas
y efectivas se hagan primero:

\begin{table}[H]
\centering
\caption{Orden sugerido de corrección.}
\label{tab:plan-correccion}

\begin{tabularx}{\textwidth}{c X l}
\toprule
\textbf{Orden} &
\textbf{Acción} &
\textbf{Hallazgo} \\
\midrule
1 & Actualizar \texttt{body-parser} a 1.20.9 (un cambio de dependencia). &
002 \\

2 & Actualizar la imagen de Nginx, fijarla por \textit{digest} y ocultar su
versión. & 003 \\

3 & Eliminar \texttt{ports: "3001:3001"} del backend, para que el acceso pase
siempre por Nginx. & S2 \\

4 & Respuestas de error genéricas y sin \texttt{X-Powered-By}. & 007 \\

5 & Restringir CORS y añadir cabeceras de seguridad en Nginx. & 004, 005 \\

6 & Fijar las acciones de GitHub por \textit{hash}. & 006 \\

7 & Autenticación y autorización en la API (cambio de mayor alcance). &
001 \\
\bottomrule
\end{tabularx}
\end{table}

\subsection{Verificación de las correcciones}

Cada corrección se validará volviendo a ejecutar el \textit{pipeline}. Un
hallazgo resuelto debe desaparecer de su reporte (Trivy, ZAP o Semgrep) y, en
la siguiente importación, DefectDojo lo marcará como cerrado. Los hallazgos
que se mantengan seguirán activos y se justificará su aceptación o su
aplazamiento. Para los que hoy no detecta ninguna herramienta, como la
ausencia de límites de tasa, se propone añadir pruebas propias o reglas
personalizadas de Semgrep.
